\documentclass[../DoAn.tex]{subfiles}
\begin{document}

\begin{center}
    \Large{\textbf{ABSTRACT}}\\
\end{center}
\vspace{1cm}
\hspace*{13pt} In the context of growing remote teamwork and Agile software development, the demand for an intuitive, real-time, and collaborative task management platform has become increasingly essential. While many current project management systems support the Kanban model, task assignment, and progress tracking, they still present several limitations, such as non-instantaneous synchronization of changes among team members, difficulty in visualizing dependencies between tasks, and a lack of automated features in task creation, description, and information synthesis.

Based on these practical needs, this thesis presents the development of HustFlow, a real-time collaborative project management platform based on the Kanban board model, integrated with artificial intelligence to assist with common management tasks. This approach is chosen because it preserves the simplicity and intuitiveness of the Kanban model, while simultaneously enhancing team collaboration and reducing manual operations in project management.

The HustFlow system is designed with core functionalities including the management of workspaces, boards, lists, and cards; member authorization; comments; activity logs; and the management of labels, checklists, attachments, and card dependencies. In addition, the platform integrates a real-time synchronization mechanism to update board modifications for active members without requiring a page refresh. Furthermore, an AI assistant is incorporated to support card creation from natural language text, standardization of task descriptions, and report generation based on board activity data.

The primary contribution of this thesis is the design and implementation of a practical task management platform that integrates the Kanban model, real-time collaboration, task dependency management, and intelligent AI assistance. Experimental results demonstrate that the system satisfies the primary functional requirements, operates stably within the testing scope, and is highly extensible to serve software development teams, academic groups, or small-to-medium organizations.

\end{document}